\documentclass[main.tex]{subfiles}

\begin{document}

\section{Heater Box Overview}

\subsection{Purpose}

The Heater Control Box is an introductory process control trainer designed to teach PLC-based temperature regulation. It provides a self-contained thermal plant with all necessary actuators, sensors, and operator interface elements wired to the PLC through the DB37 passthrough board. Students use this box to learn digital and analog I/O configuration, PID loop tuning, and alarm handling in a safe, low-power environment.

\begin{figure}[H]
    \centering
    \includegraphics[width=\textwidth]{"images/Heater Box Schematic.png"}
    \caption{Heater Control Box schematic.}
    \label{fig:heaterbox-schematic}
\end{figure}

\subsection{Theory of Operation}

The control loop operates as follows:

\begin{enumerate}
    \item The PLC energizes a digital output to close the \textbf{heater relay}, which switches current through a \textbf{150\,\si{\ohm} resistive heater element}.
    \item The heater (resistor) wams and is monitored by a \textbf{Type K thermocouple}.
    \item The thermocouple signal is conditioned by a \textbf{ProSense XTD2\_TT} temperature transmitter, which converts it to a 4--20\,mA current loop signal proportional to temperature.
    \item The 4--20\,mA signal feeds an analog input channel on the PLC (5069-IY4), where it is scaled to engineering units.
    \item The PLC executes a control algorithm (on/off, P, PI, or PID) and modulates the heater relay output accordingly.
    \item A \textbf{fan}, controllable via a separate PLC digital output, can be activated to introduce a thermal disturbance or to accelerate cooling, providing students with a means to test disturbance rejection.
\end{enumerate}

\subsection{Operator Interface Elements}

The front panel provides the following elements for student interaction and process visualization:

\begin{table}[H]
\centering
\caption{Heater Box Front Panel Elements}
\label{tab:front-panel}
\begin{tabularx}{\textwidth}{l l X}
\hline
\textbf{Element} & \textbf{Quantity} & \textbf{Function} \\
\hline
Indicator LEDs & 5 (2$\times$ Green, Yellow, 2$\times$ Red) & Status indication driven by PLC digital outputs (e.g., heater on, alarm, run/stop) \\
Pushbuttons & 2 & Operator inputs to PLC digital inputs (e.g., start/stop, acknowledge) \\
Potentiometer & 1 & Analog setpoint input (voltage divider providing 0--10\,V to PLC analog input) \\
Fan & 1 & Thermal disturbance element, driven by PLC digital output \\
Analog Voltmeter & 1 & Visual display of a PLC analog output (e.g., process variable or control output) \\
\hline
\end{tabularx}
\end{table}

\end{document}
